% ======================================================================
% SECTION 8: HR 4506 - American Leadership in Medical AI and Robotics Act of 2026
% New Federal Authorizing Statute Adapting FDORA Section 3209 and
% 21st Century Cures Act 2016.
% ======================================================================

[PROCESSING INSTRUCTION (HR 4506 - American Leadership in Medical AI
and Robotics Act of 2026, 7 to 10 paragraphs structured into the five
legislative-act subsections):

The bill is the project brief's "United States Number 1" statute. It
must explicitly state in its first paragraph that it is a "New Federal
Authorizing Statute Adapting FDORA Section 3209 (PL 117-328) and the
21st Century Cures Act of 2016 (PL 114-255)" so that readers can trace
the chain of prior authority. Cite
\cite{rank4_fdora_fda_modernization_2022} and
\cite{rank1_cures_act_2016}. Maintain the respectful FDA framing
throughout: this is an authorization, not a critique.

Subsection 8.1 - Findings and Declarations: Read in full
- national-platform/README.md (the National Platform overview)
- national-platform/new_paper/main.tex and
  national-platform/new_paper/sections/executive_summary.tex,
  introduction.tex, regulatory_landscape.tex, gov_framework.tex,
  implementation_strategy.tex, financial_analysis.tex,
  site_establishment.tex, and conclusion.tex (read all eight)
- patients/paper/Deep-Research-3/part2_ai_robotics_legislation.md (the
  legislative-vehicle taxonomy with the federal authorizing statute
  ranked first)

Compose 5 to 7 enumerated findings citing
\cite{kawchak_2026_19244918}, \cite{kawchak_2026_18445179} (replace
with main_repo cite key in the .bib),
\cite{rank4_fdora_fda_modernization_2022}, \cite{fda_rtct_press_2026},
and \cite{fda_oce_advancing_oncology_dct_2024}. Each finding must
state explicitly that the United States must remain Number 1 in the
world for patient benefit during the medical AI and robotics
revolution. At least one finding must reference the FDA's own April
2026 RTCT announcement as evidence of the FDA's leadership posture
that this bill formalizes into statute.

Subsection 8.2 - Definitions:
- "Medical AI and robotics revolution" means the period from
  approximately 2024 through 2030 during which Physical AI systems
  and large-context AI agents (e.g., Claude Code Opus 4.7 with 1M
  token context, see \cite{claude_opus_47}) are entering routine
  clinical trial deployment. Reference the four-simulation evidence
  base at new-trial/national-24-7-trial/paper/full-paper/sections/results.tex.
- "American leadership" means a documented United States lead in:
  (a) the number of authorized Physical AI oncology trial sites,
  (b) the number of patient-self-selected trial bookings per year,
  (c) the cumulative documented error-rate reduction across CMS-
  billable sites, (d) the number of patients whose data is governed
  under HR 4507 self-custody, and (e) the number of FDA RTCT pilot
  trials with patient-sponsor direct channels under HR 4505.
- "Federal investment" means dedicated appropriations to FDA, NIH
  (PAR-25-170 and successor programs, see
  \cite{nih_par25_170_dht_endpoints}), AHRQ, and ONC for the seven-
  year period 2026 through 2032.

Subsection 8.3 - American Leadership Authorizations (operative core):
- Authorization for the FDA RTCT pilot program to expand from the
  initial two trials (TRAVERSE, STREAM-SCLC) to at least 50 trials
  by 2028. Cite \cite{fda_rtct_press_2026}.
- Authorization for FDA Oncology Center of Excellence to publish
  per-state dashboards of Physical AI oncology trial site
  authorization counts, harmonized with the prior site
  documentation framework at
  new-trial/site/05-state-regulations/ca_state_regulations.tex and
  national-platform/new_trial_psl/05_title22_regulations.tex.
- Authorization for federated-learning infrastructure across CMS-
  billable sites, drawing on the federated-learning framework at
  national-platform/federated_learning/main_chunk1_preamble_intro_methods.tex,
  main_chunk2_results_architecture_pillars_examples.tex,
  main_chunk3_results_peerreview_trust_analytics.tex, and
  main_chunk4_discussion_limitations_conclusion.tex (read all four
  chunks).
- Authorization for federated investment in foundation models for
  pathology (Virchow), prediction (SCORPIO, PROGPATH), and trial
  matching (TrialGPT, PRISM, TrialMatchAI), with annual per-model
  reporting on patient-benefit metrics. Cite
  \cite{vorontsov2024virchow}, \cite{jin2024trialgpt},
  \cite{gupta2024prism}, and \cite{abdallah2026trialmatchai}.
- Authorization for a public "American Leadership Index" published
  annually with the five metrics defined in 8.2 above.

Subsection 8.4 - Implementation: 12-month deadline for the first
American Leadership Index publication; 24-month deadline for FDA
RTCT pilot expansion to 25 trials (target 50 by 2028); FDA, NIH,
AHRQ, and ONC roles. Reference
national-platform/new_template/sections/implementation_strategy.tex
for the prior implementation pattern.

Subsection 8.5 - Reporting and Enforcement: Annual American
Leadership Index report; Congressional review under FDORA
authorities; explicit alignment with the FDORA Section 3209 in silico
authorization. Cite \cite{rank4_fdora_fda_modernization_2022}.

Close the bill (1 paragraph after Subsection 8.5) with two patient-
control framings:
(i) FOR INDIVIDUAL PATIENTS: A cancer patient anywhere in the United
    States can verify, in the published American Leadership Index,
    that their state has at least one authorized Physical AI
    oncology trial site, that the per-task error rates have improved
    year-over-year, and that the FDA RTCT pilot covers their cancer
    type. The patient PAT-2026-0042 in patient-journey/ illustrates
    the per-patient outcome trajectory.
(ii) FOR BROAD ADOPTION: All 50 states must have at least one
     authorized Physical AI oncology trial site within 36 months;
     the United States retains Number 1 status in the world for
     patient benefit during the medical AI and robotics revolution
     by virtue of the per-state, per-cancer, per-site dashboard
     publication; federal appropriations sustain the FDA, NIH,
     AHRQ, and ONC infrastructure for seven years. Cite
     \cite{kawchak_2026_19244918}, \cite{kawchak_2026_19994945},
     and \cite{kawchak_2026_19396256}.]
